\section{Appendix: Measure Theoretic Probability}%
%\subsection*{Crash Course: Measure Theoretic Probability - Continued}%
\label{app:measure-theory}

This appendix provides a crash course (or rather, refresher) of concepts from measure theoretic probability.

\subsection{Why Measure Theory?}%

%
\paragraph*{Discrete and absolute continuous distributions are not general enough}%

\begin{Eg}[Simple example of a non-discrete non-absolute-continuous distribution]%
    Consider a uniformly distributed random variable on the interval $\Xcal:=[0,1]$, i.e.\ $X \sim \Ucal[0,1]$, which has probability density:
    \[ p(x) = \I_{[0,1]}(x).\]
    Consider an exact copy of $X$, which we call $Y:=X$, on $\Ycal :=[0,1]$. Now consider the joint distribution of $(X,Y)$ 
    on $\Xcal \times \Ycal = [0,1]^2$.
    Then only values on the diagonal $\Delta:=\{ (x,x) \,|\, x \in [0,1]\}$ can be realized by $(X,Y)$.
    This simple distribution on $[0,1]^2$ is not discrete (as it can attain uncountably many values), and it is also not 
    absolute continuous, since we have: $\int_\Delta \,dx\,dy = 0$, i.e.\ the (2-dimensional) area of the (1-dimensional) line is zero. 
    This implies that any density function $p$ would satisfy: $\int_\Delta p(x,y)\,dx\,dy =0$ as well. This is in contrast to the fact that a probability distribution should always be normalized:
    \[ 1 = P((X,Y) \in \Delta) = \int_\Delta p(x,y)\,dx\,dy.\]
    Note that we don't need a probability density to be able to assign probabilities to subsets $D \ins [0,1]^2$.
    We can just use  the push-forward map:
    \[(X,Y):\; [0,1] \to [0,1] \times [0,1], \quad x \mapsto (x,x).\]
    and compute:
    \[P((X,Y) \in D) = P(\{x \in [0,1]\,|\, (x,x) \in D\}),\]
    where $P$ on the right here denotes the uniform distribution on $[0,1]$.
\end{Eg}%

\begin{Not}[Unifying the notations to measure theoretic ones]%
    Let $X$ be a  random variable taking values in space $\Xcal$ and with probability distribution $P$.
    Let $F: \Xcal \to \R$ be a function. Then we will change the notations for expectation values as follows.
 \begin{enumerate}
     \item Let $X$ be a \emph{discrete} random variable with \emph{probability mass function} $p$. Then define:
 \begin{align*}
     \E[F(X)] &= \sum_{x \in \Xcal} F(x) \cdot p(x) \\
         &=: \int F(x)  \, P(dx)\\
         &=: \int F(x)  \, dP(x)\\
         &=: \int F  \, dP.
 \end{align*}
  We will consider sums to be special cases of measure integrals.
     \item Let $X$ be a \emph{absolute continuous} random variable with \emph{probability density function} $p$. Then define:
 \begin{align*}
     \E[F(X)] &= \int_\Xcal F(x) \cdot p(x)\,dx \\
         &=: \int F(x)  \, P(dx)\\
         &=: \int F(x)  \, dP(x)\\
         &=: \int F  \, dP.
 \end{align*}
 \end{enumerate}
 So both cases can be unified with the 3 commonly used notations:
 \[ \E[F(X)] = \int F\,dP = \int F(x)\,dP(x) = \int F(x)\,P(dx).  \]
 Note that in both cases we also can write: $P(A) = \int \I_A \, dP$.
\end{Not}

\begin{Exc}
    Show that the following relation holds:
    \[ \int F(x)\,P(dx) = \int z \,P^F(dz).\]
\end{Exc}

\paragraph*{Defining probability distributions on all subsets is too general}%

\begin{Rem}\label{rem-power-set}
    When we want to work with a (probability) measure $\mu$ we at least want to require that it is \emph{countably additive}, i.e.\ that for pairwise \emph{disjoint} subsets $A_n \ins \Xcal$, $n \in \N$, we have:
    \[\mu\lp\bigcup_{n \in \N} A_n\rp = \sum_{n \in \N} \mu(A_n).\]
%    where $\bigdcup$ is the symbol for disjoint union.
    We will see below that if we do not restrict the subsets $A_n$ in some way we will encounter strange behaviour.
\end{Rem}

\begin{Thm}[Vitali, non-existence of Lebesgue measure on all subsets]%
   \label{thm-non-lebesgue}
    There does NOT exist a measure $\lambda$ on $[0,1]$ such that:
    \begin{enumerate}
        \item $\lambda$ can measure every subset $A \ins [0,1]$, and:
        \item $\lambda\lp[a,b] \rp = b-a$ for all $a\le b$ with $a,b \in [0,1]$.
    \end{enumerate}
    In other words, there does NOT exist a uniform distribution on $[0,1]$ that can consistently assign values to all subsets.\\
    But: such a measure with property 2. exists on the set $\Bcal_{[0,1]}$ of all so called \emph{Borel subsets} (or even Lebesgue subsets) 
    of $[0,1]$.
    %It is called the \emph{Lebesgue measure} on $\R$ (or, more precisely, on $\Bcal_\R$). 
    Similar statements hold for higher dimensions and $\R^D$ and higher dimensional volumes.
\end{Thm}

\begin{Eg}[Vitali set]
    Consider the following equivalence relation on $[0,1]$:
    \[ r_1 \sim r_2 \qquad :\iff \qquad r_2-r_1 \in \Q.\]
    Let $[0,1]/\mathord{\sim}$ be the set of equivalence classes.
    By the axiom of choice there exists a representative system $V \ins [0,1]$ for $[0,1]/\mathord{\sim}$.
    This means that the map:
    \[ V \to [0,1]/\mathord{\sim}, \qquad v \mapsto [v],\]
    is bijective. $V$ is called \emph{Vitali set} and we claim that $V$ is not Lebesgue-measurable.\\
    For this let $q \in \Q$ and consider the subset:
    \[V_q :=  V+q := \lC v + q \,|\, v \in V \rC  \ins \R. \]
    Let $\Qcal := [-1,1] \cap \Q$.  Note that $[0,1]$ and $V_q$ are uncountable while $\Q$ and $\Qcal$ are countably infinite. 
    We then have the inclusions:
    \[ [0,1] \ins  \bigcup_{ q \in \Qcal} V_q \ins [-1,2].\]
    The right inclusion is clear as:
     \[V + [-1,1] \ins [0,1]+[-1,1] \ins [-1,2].\]
    For the left inclusion let $ x\in [0,1]$. By construction there exists a $v \in V$ such that $v \sim x$. 
    So $x - v \in \Q$. Since $x,v \in [0,1]$ we also have that $x-v \in [-1,1]$. So $q := x-v \in [-1,1] \cap \Q =\Qcal$.
    This shows that $ x \in V_q$ for a $q \in \Qcal$. Thus both inclusions are shown.\\
    If we now assumed that $V$ would be Lebesgue-measurable then every $V_q$ would be as well as a translated version of $V$.
    We then would get that: $\lambda(V_q) = \lambda(V)$ for every $q\in \Q$. So we would get:
    \[ 1 = \lambda \lp[0,1]\rp \le \lambda\lp\bigcup_{ q \in \Qcal} V_q \rp   \le \lambda \lp[-1,2]\rp =3  ,\]
    which implies:
    \[ [1,3] \ni \lambda\lp\bigcup_{ q \in \Qcal} V_q \rp = \sum_{ q\in \Qcal} \lambda(V_q) = 
    \sum_{ q\in \Qcal} \lambda(V), \]
    which is contradictory. Indeed, $\lambda(V) =0$ can be ruled out as the sum would sum up to $0 \notin [1,3]$.
    But also $\lambda(V)>0$ can be ruled out as this would sum up to $\infty \notin [1,3]$.
    So the \emph{Vitali set} $V$ can not be Lebesgue-measurable.
\end{Eg}




\begin{Thm}[Banach-Tarski paradox]
    \label{thm-banach-tarski}
    The 3-dimensional unit ball $B_1(z)=\{ x \in \R^3\,|\, \|x-z\| \le 1\}$ centered at $z \in \R^3$ can be partitioned into a finite number of disjoint sets $A_1,\dots,A_K$ (e.g.\ $K=5$) such that each can then be rotated and translated in $\R^3$ such that they form TWO 3-dimensional unit balls $B_1(y_1)$ and $B_1(y_2)$.\\
    Note that the unit balls have well-defined volume (i.e.\ 3-dimensional Lebesgue measure) and translation and rotations are very well behaved and preserve volume, while the subsets $A_k$ are very pathological (i.e.\ non-Lebesgue-measurable).
\end{Thm}


\begin{figure}[ht]
    \centering
\includegraphics[scale=0.33]{Banach-Tarski-Paradox}
\caption[Banach-Tarski paradox]{Illustration of the Banach-Tarski paradox.\footnotemark}
\label{fig:banach-tarski}
\end{figure}
\footnotetext{\href{https://en.wikipedia.org/wiki/Banach\%E2\%80\%93Tarski\_paradox
}{https://en.wikipedia.org/wiki/Banach-Tarski\_paradox} }



\textbf{$\Longrightarrow$ Measure theory is the unifying `safe space' for probability theory!}


\subsection{Core Concepts}

\begin{Motivation}
    As discussed before in remark \ref{rem-power-set}, we want to define probability measures $P$ on a space $\Wcal$. We want them to follow (at least) these rules:
    \begin{enumerate}[label=\roman*)]
        \item normalized: $P(\Wcal)=1$, $P(\emptyset)=0$.
        \item complement: $P(A^\cmpl) = 1 - P(A)$ for $A \ins \Wcal$.
        \item $\sigma$-additivity (aka countably additivity): For pairwise disjoint subsets $A_n \ins \Wcal$, $n \in \N$:
    \[P\lp\bigcup_{n \in \N} A_n\rp = \sum_{n \in \N} P(A_n).\]
    \end{enumerate}
    Such rules implicitly assume that $P$ can measure the sets $\Wcal$ and $\emptyset$; 
    and that $P$ can measure the complement $A^\cmpl$ if it can measure $A$; 
    and that $P$ can measure the (disjoint) union $\bigcup_{n \in \N} A_n$ if it can measure each of the $A_n$.\\
    As illustrated by the theorems \ref{thm-non-lebesgue} and \ref{thm-banach-tarski}, this is in general NOT possible to do for all subsets of $\Wcal$ (i.e.\ for all elements of the power set $2^\Wcal$).\\
    This problem is solved and formalized by the notion of \emph{$\sigma$-algebras} of subsets of the space $\Wcal$.
\end{Motivation}

%\subsection{Short Overview}
%- Def sigma algebras
%- Def measurable spaces
%- Def measurable maps
%- Def measures



\begin{Def}[$\sigma$-algebras]
    \label{def-sigma-algebra}
    Let $\Wcal$ be a set. A (non-empty) set $\Bcal \ins 2^\Wcal$ of subsets $A \ins \Wcal$ is called a 
    \emph{$\sigma$-algebra} on $\Wcal$ if it satisfies the following rules:
    \begin{enumerate}[label=\roman*)]
        \item empty set: $\emptyset \in \Bcal$,
        \item complement: If $A \in \Bcal$ then also: $A^\cmpl:= \Wcal\sm A \in \Bcal$,
        \item countable union: If $A_n \in \Bcal$ for all $n \in \N$ then also: $\bigcup_{n \in \N} A_n \in \Bcal$.
    \end{enumerate}
\end{Def}

\begin{Def}[Measurable spaces]
    \label{def-measurable-space}
    A tuple $(\Wcal,\Bcal)$ of a set $\Wcal$ and a $\sigma$-algebra $\Bcal$ on $\Wcal$ is called \emph{measurable space}.
\end{Def}

\begin{Rem}[Abuse of notation]
 By abuse of notation we often just call $\Wcal$ a measurable space by implicitly assuming that it is endowed with a fixed $\sigma$-algebra, which we will indicate by $\Bcal_\Wcal$ or $\Bcal(\Wcal)$ if needed. We will also just call a subsets $A \ins \Wcal$ \emph{measurable} when we actually mean that $A \in \Bcal_\Wcal$.
\end{Rem}

\begin{Def}[Measures]
    \label{def-measure}
    Let $(\Wcal,\Bcal)$ be a measurable space.
    A \emph{measure} $\mu$ on $(\Wcal,\Bcal)$ - by definition - is a mapping:
    \[\mu: \; \Bcal \to \R \cup \{\infty\},\quad D \mapsto \mu(D), \]
    such that:
    \begin{enumerate}[label=\roman*)]
        \item non-negative: $\forall A \in \Bcal$: $\mu(A) \in [0,\infty]$,
        \item empty set: $\mu(\emptyset) =0$,
        \item countably additive (aka $\sigma$-additive): for all sequences $A_n \in \Bcal$, $n \in \N$, with $A_i \cap A_j = \emptyset$ for all $i \neq j$, we have:
    \[\mu\lp\bigcup_{n \in \N} A_n\rp = \sum_{n \in \N} \mu(A_n).\]
    \end{enumerate}
    \end{Def}

\begin{Def}[Probability/finite/$\sigma$-finite measures]
    A measure $\mu$ on $(\Wcal,\Bcal)$ is called:
    \begin{enumerate}
        \item \emph{probability measure} if $\mu(\Wcal)=1$.
        \item \emph{finite measure} if $\mu(\Wcal) < \infty$.
        \item \emph{$\sigma$-finite measure} if there are $D_n \in \Bcal$, $n \in \N$, with $\mu(D_n) < \infty$ and $\Wcal = \bigcup_{n \in \N} D_n$.
    \end{enumerate}
\end{Def}

\begin{Def}[Measure spaces/probability spaces]
    A triple $(\Wcal, \Bcal, \mu)$ consisting of a measurable space $(\Wcal,\Bcal)$ and a measure $\mu$ on $(\Wcal,\Bcal)$ is called \emph{measure space}
    (and \emph{probability space} if $\mu$ is a probability measure).\\
    Again, by abuse of notation, we often omit the $\sigma$-algebra in the notation and call $(\Wcal,\mu)$ a measure space, probability space, resp.
\end{Def}

\begin{Def}[Measurable mappings]
%    \label{def-measurable-map}
    Let $(\Wcal,\Bcal_\Wcal)$ and $(\Zcal,\Bcal_\Zcal)$ be two measurable spaces and $f: \Wcal \to \Zcal$ be a mapping.
    We call $f$ a \emph{$\Bcal_\Wcal$-$\Bcal_\Zcal$-measurable mapping} (or just \emph{measurable} for short) if
    for all $B \in \Bcal_\Zcal$ the pre-image $f^{-1}(B)$ is an element of $\Bcal_\Wcal$. In formulas:
    \[\forall B \in \Bcal_\Zcal:\; f^{-1}(B) \in \Bcal_\Wcal.\]
    Remember the definition of pre-image: $f^{-1}(B) := \{ w\in \Wcal\,|\, f(w) \in B \}$.
\end{Def}



\begin{Def}[Push-forward measure]
    Let  $X:\; (\Wcal,\Bcal_\Wcal) \to (\Xcal,\Bcal_\Xcal)$ be measurable and $\mu$ a measure on $(\Wcal,\Bcal_\Wcal)$.
    Then we define the \emph{push-forward measure} (aka \emph{image measure}) of $\mu$ via:
    \[ (X_*\mu)(A):= \mu^X(A) := \mu_X(A) := \mu(X)(A):= \mu(X \in A):= \mu(X^{-1}(A))\]
    for all $A \in \Bcal_\Xcal$.
    If $\mu$ is a probability distribution then the push-forward measure $\mu(X)$ is also called the (distributional) \emph{law of $X$}.
\end{Def}


\begin{Def}[Random variables]
    A measurable mapping: \[X:\; (\Wcal,\Bcal_\Wcal,P) \to (\Xcal,\Bcal_\Xcal)\] 
    that starts from a probability space is also called \emph{random variable}.\\
    The main point is that the map $X$ comes with its own distribution $P^X$.
    We often just say: ``Let $X$ be a random variable with distribution $P^X = \cdots$'',
    where $P^X$ is then specified, e.g.\ to be a Gaussian or a categorical distribution, etc.
\end{Def}


\begin{Def}[Null sets]
    Let $(\Wcal,\Bcal,\mu)$ be a measure space.
    A subset $M \ins \Wcal$ is called \emph{$\mu$-null} or \emph{$\mu$-zero set} if there exists a set $N \in \Bcal$ with $M \ins N$ and $\mu(N)=0$.
\end{Def}

\begin{Def}[Almost surely/almost all]
    Let $(\Xcal,\Bcal,\mu)$ be a measure space and $f,g:\,\Xcal \to \Zcal$ a measurable map.
    We write $f=_\mu g$ or say $f=g$ \emph{$\mu$-almost-surely} (a.s.) or $f(x)=g(x)$ for \emph{$\mu$-almost-all} $x \in \Xcal$ 
    if:  
    \[\{ x \in \Xcal\,|\,f(x)\neq g(x) \} \quad \text{ is a $\mu$-null set.} \]
    Similarly, for $f\le_\mu g$, etc.. \\
    More generally,  we say that a condition $C$ about points $x \in \Xcal$ holds \emph{$\mu$-almost-surely} or for \emph{$\mu$-almost-all} $x \in \Xcal$ if the set of points where the condition does not hold is $\mu$-null, i.e.:
    \[\{ x \in \Xcal\,|\,\neg C(x) \} \quad \text{ is a $\mu$-null set}. \]
\end{Def}

\subsection{Default Choices for Sigma-Algebras}

%\ref{exercise-sigma-algebra}

In this subsection we want to highlight what kind of default $\sigma$-algebras we will assume 
on different types of spaces and on spaces constructed from others.

\begin{Rem}[Discrete spaces]
    If $\Wcal$ is countable (i.e.\ either finite or countably infinite, e.g.\ like $\Z$, $\Q$ or  $\N$ or $\{1,\dots,N\}$) 
    then we will always implicitly assume that $\Wcal$ is endowed with the power set $\sigma$-algebra: $\Bcal_\Wcal = 2^\Wcal$ (unless stated otherwise).
\end{Rem}


\begin{Def}[$\sigma$-algebra generated by a set of subsets]
    \label{def-sigma-algebra-gen}
    Let $\Wcal$ be a set and $\Acal \ins 2^\Wcal$ be any non-empty set of subsets of $\Wcal$.
    Then we can define the \emph{$\sigma$-algebra generated by $\Acal$}:
    \[ \sigma(\Acal) := \bigcap_{\substack{ \Acal \ins \Bcal \\\Bcal\; \sigma\text{-algebra on }\Wcal }} \Bcal, \]
    as the intersection of all $\sigma$-algebras $\Bcal$ on $\Wcal$ that contain $\Acal$. 
    Note that the set $\sigma(\Acal)$ really is a well-defined $\sigma$-algebra on $\Wcal$. % (see exercise \ref{exercise-sigma-algebra}).\\
    $\sigma(\Acal)$ is thus - by definition - the smallest $\sigma$-algebra on $\Wcal$ that contains $\Acal$.
\end{Def}


\begin{Def}[Borel $\sigma$-algebra on topological spaces]
    Let $(\Wcal,\Ocal)$ be a topological space with set of open subsets $\Ocal$ then the \emph{Borel $\sigma$-algebra} of $(\Wcal,\Ocal)$ is 
    defined as the smallest $\sigma$-algebra that contains all open (and thus also all closed) subsets:
    \[ \Bcal_{(\Wcal,\Ocal)} := \sigma(\Ocal).\]
    We will always implicitly assume that every topological space is endowed with its Borel $\sigma$-algebra (unless stated otherwise).
\end{Def}

\begin{Rem}
    Caution: Other choices of $\sigma$-algebras for topological spaces used in the literature are the Baire $\sigma$-algebra, which is generated by the zero sets of all continuous functions, or the $\sigma$-algebra generated only by its closed (countably) compact sets, or the $\sigma$-algebra of all (Radon-)universally measurable subsets.
\end{Rem}

\begin{Lem}[Borel $\sigma$-algebra on $\R^D$]    
  %  Consider the Euclidean space $\R^D$. Then we define the \emph{Borel $\sigma$-algebra} $\Bcal_{\R^D}$ 
   % on $\R^D$ as the $\sigma$-algebra generated by all the closed compact cubes/intervals:
    The Borel $\sigma$-algebra of $\R^D$ is generated by the cubes:
\[ \Bcal_{\R^D} = \sigma\lp\lC [a_1,b_1] \times \cdots \times [a_D,b_D]\;|\; a_d, b_d \in \Q, a_d \le b_d, d=1,\dots,D    \rC\rp.  \]
%We will always implicitly assume $\R^D$ to be endowed with this Borel $\sigma$-algebra $\Bcal_{\R^D}$ (unless stated otherwise).
\end{Lem}

\begin{DefLem}[$\sigma$-algebras induced by mappings]
    \label{sigma-algebra-map}
    Let $f: \Wcal \to \Zcal$ be any mapping.
    \begin{enumerate}
        \item Let $\Bcal_\Zcal$ be a $\sigma$-algebra on $\Zcal$. Then the \emph{pull-back $\sigma$-algebra} defined via:
            \[ f^*\Bcal_\Zcal := \lC f^{-1}(C)\,|\, C \in \Bcal_\Zcal   \rC \]
            is the smallest $\sigma$-algebra $\Bcal_\Wcal$ that makes $f$ $\Bcal_\Wcal$-$\Bcal_\Zcal$-measurable.
        \item Let $\Bcal_\Wcal$ be a $\sigma$-algebra on $\Wcal$. Then the \emph{push-forward $\sigma$-algebra} defined via:
            \[ f_*\Bcal_\Wcal := \lC C \ins \Zcal\,|\, f^{-1}(C) \in \Bcal_\Wcal   \rC \]
            is the biggest $\sigma$-algebra $\Bcal_\Zcal$ that makes $f$ $\Bcal_\Wcal$-$\Bcal_\Zcal$-measurable.
    \end{enumerate}
\end{DefLem}


\begin{Def}[Product $\sigma$-algebra]
    Let $(\Xcal_i,\Bcal_i)$ be measurable spaces, $i \in I$.
    Then the product space $\prod_{i \in I} \Xcal_i$ is endowed with the smallest $\sigma$-algebra
    such that for every $j\in I$ the projection map:
    \[ \pr_j: \prod_{i \in I} \Xcal_i \to \Xcal_j,\qquad (x_i)_{i \in I} \mapsto x_j, \]
    is measurable. We use the symbols $\bigotimes_{i \in I} \Bcal_i$ for this product $\sigma$-algebra.
    In symbols:
    \[ \bigotimes_{i \in I} \Bcal_i := \sigma\lp \bigcup_{i \in I} \pr_i^*\Bcal_i\rp. \]
  We will always implicitly assume that every product space is endowed with this product $\sigma$-algebra (unless stated otherwise).
\end{Def}

%\begin{Rem}
%    \begin{enumerate}
%        \item  Note that not every $C \in \Bcal_\Xcal \otimes \Bcal_\Ycal$ is of such product form $A \times B$.
%        \item  It is possible to define the product $\sigma$-algebra: $\bigotimes_{i \in I}\Bcal_{\Xcal_i}$ on the product $\prod_{i \in I} \Xcal_i$ 
%            of arbitrary measurable spaces $(\Xcal_i, \Bcal_{\Xcal_i})$, $i \in I$, similarly (with a bit of care for infinite index sets $I$).
%    \end{enumerate}
%\end{Rem}


%\ref{exercise-Borel-sigma-algebra}


\begin{Def}[Subspace $\sigma$-algebra]
    Let $(\Wcal,\Bcal)$ be a measurable space and $\Zcal \ins \Wcal$ be a subset.
    Then the \emph{subspace $\sigma$-algebra} $\Bcal_{|\Zcal}$ on $\Zcal$ is the smallest $\sigma$-algebra that makes the inclusion map $\Zcal \to \Wcal$ measurable. More concretely:
    \[\Bcal_{|\Zcal} := \lC B \cap \Zcal \;|\; B \in \Bcal   \rC.\]
    We will always assume that subsets are endowed with the subspace $\sigma$-algebra (unless it is ambiguous or stated otherwise).
\end{Def}

\begin{Def}[Disjoint union $\sigma$-algebra]
 Let $(\Xcal_i,\Bcal_i)$ be measurable spaces, $i \in I$, considered to be pairwise disjoint.
 Then the \emph{disjoint union $\sigma$-algebra} on the disjoint union 
 $ \coprod_{ i \in I} \Xcal_i$ 
 is the biggest $\sigma$-algebra $\Bcal_\sqcup $ such that all inclusion maps $\Xcal_i \to \coprod_{i \in I} \Xcal_i$ are measurable.
 In symbols:
 \[ \Bcal_\sqcup := \lC E \ins {\coprod_{ i \in I} \Xcal_i} \,\bigg|\,  \forall i \in I:\, E \cap \Xcal_i \in \Bcal_i\rC. \]
\end{Def}

\begin{Def}[$\sigma$-algebra on the space of all probability measures]
    \label{def-sigma-algebra-prob-space}
    Let $(\Wcal,\Bcal_\Wcal)$ be a measurable space. We denote the space of all probability measures on  $(\Wcal,\Bcal_\Wcal)$ by:
    \[ \Pcal(\Wcal) := \lC P \;|\; P \text{ is probability measure on } (\Wcal,\Bcal_\Wcal) \rC.\]
    We endow $\Pcal(\Wcal)$ with the smallest $\sigma$-algebra $\Bcal_{\Pcal(\Wcal)}$ such that all evaluation maps:
    \[  \ev_D:\; \Pcal(\Wcal) \to [0,1],\qquad P \mapsto P(D)  \]
    are measurable for $D \in \Bcal_\Wcal$.
    In symbols:
    \[ \Bcal_{\Pcal(\Wcal)} := \sigma\lp \bigcup_{D \in \Bcal_\Wcal} \ev_D^*\Bcal_{[0,1]} \rp.\]
    We will always assume that the space of probability measures $\Pcal(\Wcal)$ is endowed with this $\sigma$-algebra (unless stated otherwise).
\end{Def}


\begin{comment}
\begin{Def}[$\sigma$-algebra on the space of all measurable mappings]
    Let $\Wcal$ and $\Xcal$ be two measurable spaces and:
    \[ \Lcal(\Wcal,\Xcal) := \Lcal^0(\Wcal,\Xcal):=\lC f:\,\Wcal \to \Xcal \text{ measurable} \rC\]
    the space of all measurable mappings from $\Wcal$ to $\Xcal$.
    We will endow $\Lcal(\Wcal,\Xcal)$ with the smallest $\sigma$-algebra such that the evaluation maps:
    \[ \ev_w:\, \Lcal(\Wcal,\Xcal),\quad f \mapsto f(w), \]
    are measurable for all $w \in \Wcal$. In symbols:
    \[  \Bcal_{\Lcal(\Wcal,\Xcal)}:=\sigma\lp \bigcup_{w \in \Wcal} \ev^*_w\Bcal_\Xcal\rp.  \]
    Note that this $\sigma$-algebra does not use the $\sigma$-algebra $\Bcal_\Wcal$ of $\Wcal$ in its definition.
    $\Bcal_\Wcal$ is only used in the definition of $\Lcal(\Wcal,\Xcal)$ itself.
\end{Def}

\begin{Rem}[$\sigma$-algebra on the space of all Markov kernels]
    Let $\Wcal$ and $\Xcal$ be two measurable spaces and:
    \[ \Kcal(\Wcal,\Xcal) := \lC K:\,\Wcal \dshto \Xcal \text{ Markov kernel} \rC = \Lcal(\Wcal,\Pcal(\Xcal)),\]
    the space of all Markov kernels from $\Wcal$ to $\Xcal$, or equivalently, the space of all 
    measurable maps from $\Wcal$ to $\Pcal(\Xcal)$.
    From the two points above $\Kcal(\Wcal,\Xcal)$ is then endowed with the smallest $\sigma$-algebra such that all evaluation maps $\ev_w$
    are measurable for $w \in \Wcal$. Since $\ev_w$ map to $\Pcal(\Xcal)$, by definition, these are measurable if and only if 
    their composition $\ev_D \circ \ev_w$ with $\ev_D$ are measurable for all $D \in \Bcal_\Xcal$.
    So we get:
    \[ \Bcal_{\Kcal(\Wcal,\Xcal)} = \sigma\lp \bigcup_{\substack{D \in \Bcal_\Xcal\\w \in \Wcal}} 
    (\ev_D \circ \ev_w)^*\Bcal_{[0,1]}   \rp   \]
\end{Rem}

\begin{Def}[$\sigma$-algebra on the $\sigma$-algebra]
    Let $(\Xcal,\Bcal_\Xcal)$ be a measurable space.
    Then $\Bcal_\Xcal$ can be made into a measurable space itself by defining the $\sigma$-algebra $\Bcal_{\Bcal_\Xcal}$ 
    as the smallest $\sigma$-algebra such that the Dirac measures:
    \[ \delta_x:\, \Bcal_\Xcal \to \{0,1\},\quad A \mapsto \delta_x(A) = \I_A(x),\]
    are measurable. In symbols:
    \[  \Bcal_{\Bcal_\Xcal}  := \sigma\lp \bigcup_{x\in \Xcal} \delta_x^{-1}(\{1\}) \rp.\]
    This $\sigma$-algebra makes the following map measurable:
    \[ \Bcal_\Xcal \to \Lcal(\Xcal,\R),\quad A \mapsto \I_A.\]
\end{Def}
\end{comment}

\subsection{Standard Measurable Spaces}

\begin{Def}[Standard measurable space]
 A measurable space $(\Wcal,\Bcal)$ is called \emph{standard measurable space} (aka standard Borel space) if it is measurably isomorphic to either:
    \begin{enumerate}
        %\item a countable space $\Xcal$ endowed with the power set $\sigma$-algebra $2^\Xcal$ (e.g.\ $\Xcal=\{1,\dots,n\}$ or $\Xcal=\N$), or:
        \item a finite measurable space $\{1,\dots,M\}$ for some $M \in \N$ endowed with the power set $\sigma$-algebra $2^{\{1,\dots,M\}}$, or:
        \item the countably infinite space $\N$ endowed with the power set $\sigma$-algebra $2^\N$, or:
        \item the unit interval $[0,1]$ endowed with its Borel $\sigma$-algebra:
           \[\Bcal_{[0,1]} =\sigma\lp\lC [a,b]\,|\, a, b \in [0,1] \cap \Q, a \le b \rC \rp.\]
        %\item an Euclidean space $\R^D$ endowed with its Borel $\sigma$-algebra: 
        %    \[\Bcal_{\R^D} =\sigma\lp\lC [a_1,b_1]\times \cdots \times[a_D,b_D]\,|\, a_d, b_d \in \R, a_d \le b_d, d=1,\dots,D \rC\rp.\]
    \end{enumerate}
    'Measurably isomorphic' means that there is a measurable mapping that has a measurable inverse.
\end{Def}

\begin{Thm}[Kuratowski et al.]
    \begin{enumerate}
        \item Every Borel subset of any complete metric space that has a countable dense subset is a standard measurable space in its Borel $\sigma$-algebra 
    (e.g.\ $\Q^D$ is countable and dense in $\R^D$).
\item Two standard measurable spaces $\Xcal$ and $\Ycal$ are measurably isomorphic iff their cardinalities $|\Xcal|$, $|\Ycal|$ are equal 
    (e.g.\ $\R^D \cong [0,1]$).
\item Countable disjoint unions and countable direct products of standard measurable spaces are standard measurable spaces.
\item If $\Wcal$ is standard measurable space then the space of its probablity measures $\Pcal(\Wcal)$ is also a standard measurable space.
    \end{enumerate}
\end{Thm}

\begin{Eg}
    Examples of standard measurable spaces are: $\R^D$, $\Q$, $\Z$, $\N$, $\{1,\dots,M\}$,  $[0,1]$, topological manifolds, countable CW-complexes, 
    every Borel set of any separable complete metric space.
\end{Eg}



\subsection{Measure Integrals}
\label{sec:measure-integral}


The construction of the measure integral $\int f\,d\mu$ follows in several steps.

\begin{Construction}[Measure integral]
    Let $(\Xcal,\Bcal_\Xcal,\mu)$ be a measure space.
    \begin{enumerate}
        \item Indicator functions: For $A \in \Bcal_\Xcal$ put:
            \[\int \I_A\,d\mu := \mu(A).\]
        \item Simple functions: For a simple function $g:\; \Xcal \to \R$ given by: 
            \[ g(x)= \sum_{n=1}^N a_n \cdot \I_{A_n}(x), \]
            where  $A_n \in \Bcal_\Xcal$ %with $\mu(A_n) < \infty$ 
            and $a_n \in \R$, $n=1,\dots,N$, we define: 
            \[ \int g\,d\mu := \sum_{n=1}^N a_n \cdot \mu(A_n). \]
        \item Non-negative measurable functions: Let $h:\, \Xcal \to [0,\infty]$ be a non-negative measurable function then we define:
            \[ \int h\, d\mu := \sup_{0 \le g \le h} \int g \, d\mu \qquad \in [0,\infty], \]
            where the supremum is running over all non-negative simple functions $g$ that are smaller or equal to $h$.
            %are only bigger than $h$ on a $\mu$-null set.
        \item Measurable functions with well-defined integral: Let $f:\, \Xcal \to \bar\R=\R \cup \{\pm \infty\}$ be a measurable function. 
            We then can  write  $f=f_+-f_-$ with:
            \[ f_+ := \max(f,0) \ge 0, \qquad f_-:= \max(-f,0) \ge 0.\]
            If at least one of $\int f_+\,d\mu$,  $\int f_-\,d\mu$ is finite (i.e.\ $< \infty$)
            we can then define:
            \[ \int f\, d\mu := \int f_+\,d\mu - \int f_-\,d\mu \qquad \in [-\infty,\infty].\]
        \end{enumerate}
   The only case where we cannot properly define the integral is for measurable functions $f:\, \Xcal \to \bar \R$ where both integrals: 
   $\int f_+ \,d\mu = \infty$ and $\int f_-\,d\mu = \infty$ are infinite, because of the ``$\infty - \infty = ?$'' problem.

\end{Construction}

\begin{Rem}[Riemann integral vs. measure integral]
    The construction of the Riemann integral (RI) and the measure integral (MI) differ only in a few points:
    \begin{enumerate}
        \item RI uses (infinitesimal) interval length on x-axis, while MI uses the measure content
            (which is also the interval length in case of the Lebesgue measure).
        \item RI decomposes the x-axis, while MI decomposes the y-axis.
        \item RI uses limits for the integration boundaries to integrate to infinity (if convergent), 
                while MI takes difference of integrals (to infinity) of $f_+$ and $f_-$ (if difference well-defined).
        \item RI integrates in direction $a$ to $b$, while MI integrates interval $[a,b]$ in an undirected fashion.
        \item If a function is Riemann integrable (RI) (e.g. continuous) on interval $[a,b]$ 
            then it is also Lebesgue integrable (MI) with the same integral value.
    \end{enumerate}
\end{Rem}

\begin{Def}[Integrable functions]
   Let $(\Xcal,\Bcal,\mu)$ be a measure space. A measurable function $f:\,\Xcal \to \R$ is called \emph{$\mu$-integrable} if:
    \[  \int |f|\, d\mu < \infty.\]
\end{Def}

\begin{Thm}[Properties of the integral]
       Let $(\Xcal,\Bcal,\mu)$ be a measure space and $f,g:\,\Xcal \to \R$ $\mu$-integrable measurable functions.
       \begin{enumerate}
           \item For $A \in \Bcal$ we have: $\int \I_A\,d\mu = \mu(A)$.
           \item Linearity: %If $f$ and $g$ are $\mu$-integrable 
               If $a,b \in \R$ then $a\cdot f+b\cdot g$ is also $\mu$-integrable and we have:
               \[ \int (a\cdot f+b\cdot g)\,d\mu = a\cdot \int f\, d\mu + b \cdot \int g\,d\mu.\]
           %\item If $f$ is $\mu$-integrable and $ c\in \R$ then $c\cdot f$ is also $\mu$-integrable and we have:
            %   \[ \int (c\cdot f)\,d\mu = c\cdot \int f\, d\mu.\]
           \item Triangle inequality: %If $f$ is $\mu$-integrable then: 
               \[\left|\int f\,d\mu\right| \le \int |f|\,d\mu < \infty.\]
           \item %If $f$ is $\mu$-integrable and 
               If $f \ge_\mu 0$ then: $\int f\,d\mu \ge 0$, with equality iff $f =_\mu 0$.
           \item Monotonicity:  %If $f$ and $g$ are $\mu$-integrable and 
               If $f \ge_\mu g$ then: $\int f\,d\mu \ge \int g\,d\mu$, with equality iff $f=_\mu g$.
           \item If $\int f\,d\mu <\infty$ then $f<_\mu \infty$.
          % \item Monotone convergence: If $f_{n+1} \ge_\mu f_n \ge_\mu 0$ for all $n\in\N$ and $f := \limsup_{n\in\N} f_n$ then: 
           %    \[ \int f_n\, d\mu \stackrel{ n \to \infty}{\longrightarrow} \int f\,d\mu. \]
          % \item Dominated convergence: If $g$ is a non-negative $\mu$-integrable
           \item The measure integral satisfies monotone convergence, dominated convergence, Fubini theorems, etc.  (see literature).
       \end{enumerate}
       Note, we use $=_\mu$ and $\ge_\mu$ to indicate that this property is (only) allowed to fail on a $\mu$-null set.
\end{Thm}


\begin{Def}[Expectation value]
    Let $(\Wcal,\Bcal,P)$ be a probability space and $X:\, \Wcal \to \R$ be a measurable function with well-defined integral.
    Then its \emph{expectation value} (w.r.t.\ $P$) is defined to be:
    \[ \E[X]: = \int X\,dP. \]
\end{Def}

\begin{Eg} Let $\Xcal$ be a measurable space and $f:\,\Xcal \to \R$ a measurable function and $\Wcal \ins \Xcal$ a countable subset.
    \begin{enumerate}
      %  \item \emph{Counting measure}. The counting measure $\#_\Wcal$ w.r.t.\ $\Wcal$ is defined on measurable $A \ins \Xcal$:
      %      \[  \#_\Wcal(A) := \sum_{w \in \Wcal} \delta_w(A) = \sum_{w \in \Wcal}\I_A(w).  \]
        \item \emph{Dirac measure}. Let $w \in \Xcal$ be a point. We define the \emph{Dirac measure} $\delta_{w}$ centered at $w$ via:
            \[ \delta_{w}(A) := \I_A(w),\]
            for all measurable $A \ins \Xcal$. Furthermore, we have:
            \[ \E[f] =  \int f(x)\,\delta_{w}(dx)   = f(w).\]
            This holds because: $f(x) = f(w)$ for $\delta_{w}$-almost-all $x \in \Xcal$.
            Let's prove the right equality more formally:
            \begin{proof}
            Consider:
            \[B:= f^{-1}(f(w))=\{x \in \Xcal\,|\,f(x)=f(w) \} \ni w.\]
            Since $f$ is measurable and $\{f(w)\} \in \Bcal_{\R}$ we also have $B \in \Bcal_\Xcal$.
            We then have the decomposition:
            \begin{align*}
                f(x) &= f(x) \cdot \I_B(x) + f(x) \cdot \I_{B^\cmpl}(x) \\
                &= f(w) \cdot \I_B(x) + f(x) \cdot \I_{B^\cmpl}(x).
        \end{align*}
        Since $w \notin B^\cmpl$ we get $\delta_{w}(B^\cmpl) =0$ and thus $\int f(x) \cdot \I_{B^\cmpl}(x)\, \delta_{w}(dx) = 0$.
        Together we get:
        \begin{align*}
            \int f(x)\,\delta_{w}(dx) & =  \int \lp f(w) \cdot \I_B(x) + f(x) \cdot \I_{B^\cmpl}(x) \rp \, \delta_{w}(dx) \\
            & =  f(w) \cdot \int  \I_B(x)\,\delta_{w}(dx) + \underbrace{\int  f(x) \cdot \I_{B^\cmpl}(x)  \, \delta_{w}(dx)}_{=0} \\
            & = f(w) \cdot \delta_{w}(B) \\
            & = f(w).
        \end{align*}
            \end{proof}
        \item \emph{Discrete distributions}. Consider a discrete probability distribution $P$ supported on the countable subset $\Wcal \ins \Xcal$. 
            Let $p$ be its mass function.
            %Let $\#_\Wcal$ be the counting measure w.r.t.\ $\Wcal$.
            %Then the mass function of $P$ is the density of $P$ w.r.t.\ $\#_\Wcal$:
            %\[\]
            %centered at the points $x_n \in \Xcal$, $n \in \N$, with corresponding probability masses $p(x_n)$.
            We then can write the corresponding probability measure $P$ on $\Xcal$ as: % evaluated at measurable $A \ins \Xcal$ as:
            \[  P = \sum_{w \in \Wcal} p(w) \cdot \delta_{w}.\]
            For measurable $A \ins \Xcal$ we then have:
             \[ P(A) = \sum_{w \in \Wcal} p(w) \cdot \delta_{w}(A) = \sum_{w \in \Wcal \cap A} p(w). \]
            Furthermore, we get:
            \[ \E[f] = \int f(x)\, P(dx) = \int f(x)\,\sum_{w \in \Wcal} p(w) \cdot \delta_{w}(dx) = \sum_{w \in \Wcal} f(w) \cdot p(w). \]
    \end{enumerate}
\end{Eg}






\subsection{Densities/Derivatives}

\begin{Def}
    Let $(\Xcal,\Bcal)$ be a measure space and $\mu$, $\nu$ two measures on it.
    We say that $\nu$ has a \emph{density} w.r.t.\ $\mu$ if there exists a non-negative measurable function $f: \Xcal \to [0,\infty]$ such that for all $A \in \Bcal$:
    \[ \nu(A) = \int \I_A\cdot f \, d\mu =:\int_A f \, d\mu  .\]
    Such a density does not always exist.
    If a density exists then it is essentially unique, in the sense that two such densities would only differ on a $\mu$-null set.
    We often use the notation: $f=\frac{d\nu}{d\mu}$ and call it `the' \emph{density} or \emph{(Radon-Nikod\'ym) derivative} of $\nu$ w.r.t.\ $\mu$.
\end{Def}

\begin{Prp}
    Let $(\Xcal,\Bcal)$ be a measure space and $\mu$, $\nu$, $\kappa$ three measures on it and $g:\,\Xcal \to \bar\R$ is either  a $\nu$-integrable or non-negative measurable function. 
    \begin{enumerate}
        \item If $\nu$ has a density w.r.t.\ $\mu$ then we have:
                \[ \int g \,d\nu = \int g \cdot \frac{d\nu}{d\mu}\,d\mu.   \]
            \item (Conic) linearity: If $\kappa$ has a density w.r.t.\ $\mu$ and $\nu$ has a density w.r.t.\ $\mu$ and $a,b \ge0$ then $a \cdot \kappa+b \cdot \nu$ has a density w.r.t.\ $\mu$ and we have:
            \[ \frac{d(a \cdot \kappa+b\cdot \nu)}{d\mu}(x) = a\cdot\frac{d\kappa}{d\mu}(x) + b\cdot\frac{d\nu}{d\mu}(x) \]
             for $\mu$-almost-all $x \in \Xcal$.
      \item Chain rule: 
    If $\nu$ has a density w.r.t.\ $\mu$ and $\mu$ has a density w.r.t.\ $\kappa$ then also $\nu$ has a density w.r.t.\ $\kappa$ and we have:
    \[ \frac{d\nu}{d\kappa}(x) = \frac{d\nu}{d\mu}(x) \cdot \frac{d\mu}{d\kappa}(x) \]
    for $\kappa$-almost-all $x \in \Xcal$.
        \item Inverse:    If $\nu$ has a density w.r.t.\ $\mu$ and $\mu$ has a density w.r.t.\ $\nu$ then we have:
            \[ \frac{d\nu}{d\mu}(x) = \lp\frac{d\mu}{d\nu}(x) \rp^{-1}    \]
            for $\mu$-almost-all $x \in \Xcal$. We can make in this context the (somewhat arbitrary) choice to put: $0^{-1}:=\infty$.
    \end{enumerate}
\end{Prp}

\begin{Def}[Absolute continuity]
    Let $\mu$, $\nu$ be two measures on a measurable space $(\Xcal,\Bcal)$.
    We say that $\nu$ is \emph{absolute continuous} w.r.t.\ $\mu$, in symbols:
    \[ \nu \ll \mu,\]
    if for every $A \in \Bcal$ with $\mu(A)=0$ also $\nu(A)=0$ holds, in short, if:
    \[ \mu(A)=0 \quad \implies \quad \nu(A)=0.\]
\end{Def}

\begin{Thm}[Radon-Nikod\'ym, see \cite{Kle20} Cor.\ 7.34]
    \label{thm:radon-nikodym}
    Let $(\Xcal,\Bcal, \mu)$ be a $\sigma$-finite measure space and $\nu$ another measure on $(\Xcal,\Bcal)$.
    Then the following two statements are equivalent:
    \begin{enumerate}
        \item $\nu$ has a density w.r.t.\ $\mu$.
        \item $\nu$ is absolute continuous w.r.t.\ $\mu$.
    \end{enumerate}
\end{Thm}




\begin{Thm}[Besicovitch density theorem, \cite{Fremlin} 472D]
    \label{besicovitch-density}
Let $\mu$ be a Radon measure on $\R^D$ (e.g.\ any finite or probability measure or the Lebesgue measure, see \ref{thm-lebesgue-measure}) and $f:\,\R^D \to \bar\R$ be any (locally) $\mu$-integrable function.
    Then we have for $\mu$-almost-all $x \in \R^D$:
    \begin{enumerate}
        \item $\displaystyle \lim_{\varepsilon \to 0} \frac{1}{\mu( B_\varepsilon(x))} \int_{B_\varepsilon(x)}f(z)\,\mu(dz) = f(x).$   
        \item $\displaystyle \lim_{\varepsilon \to 0} \frac{1}{\mu( B_\varepsilon(x))} \int_{B_\varepsilon(x)}|f(z) - f(x)|\,\mu(dz) = 0$. 
    \end{enumerate}
    Here $B_\varepsilon(x)$ denote the closed balls of radius $\varepsilon>0$ centered at $x$ (in Euclidean norm).\\
    The above, in particular, holds for the density $f=\frac{d\nu}{d\mu}$ of another measure $\nu$ w.r.t.\ $\mu$:
    \[  \lim_{\varepsilon \to 0} \frac{\nu( B_\varepsilon(x))} {\mu( B_\varepsilon(x))}  = \frac{d\nu}{d\mu}(x), \]
    for $\mu$-almost-all $x \in \R^D$.
\end{Thm}

\subsection{Conditional Expectation}

You may be familiar with the conditional expectation for discrete random variables $X,Y$:
\begin{align*} 
    \E[X|Y=y] &= \sum_{x\in\Xcal} x\cdot P(X=x|Y=y) 
&= \sum_{x \in \Xcal} x \cdot \frac{P(X=x,Y=y)}{P(Y=y)} \\&= \frac{\sum_{x \in \Xcal} x \cdot P(X=x,Y=y)}{P(Y=y)},
\end{align*}
and for real-valued random variables $X,Y$ with positive and continuous joint density $p(x,y)$: 
\[\E[X| Y=y] = \int_\Xcal x \cdot p(x| Y=y) \,dx = \int_\Xcal x \cdot \frac{p(x,y)}{p(y)} \,dx = \frac{\int_\Xcal x \cdot p(x,y)\,dx}{p(y)}.\]
The following construction generalizes this notion:
\begin{Def}[Conditional expectation]
    Let $(\Wcal,P)$ be a probability space
    and $X:\, \Wcal \to \R$, $Y:\,\Wcal \to \Ycal$ be two random variables with 
    either $\E[|X|]< \infty$ or $X \ge 0$ a.s.
    %$\min_\pm \E[X_\pm] < \infty$.
    %$\min \lp \E[X_+], \E[X_-]\rp<\infty$.
\begin{enumerate}
         \item     The \emph{conditional expectation}  of $X$ given $Y=y$ is defined via:
     \[ \E[X|Y=y] := \E[X_+|Y=y] - \E[X_-|Y=y] \qquad \in \bar \R, \]
    % (if not both terms are $\infty$),
     where $X_\pm:=\max(\pm X,0) \ge 0$ and:
     \[\E[X_\pm|Y=y] := \frac{dE_\pm}{dP^Y}(y),\]
     is the Radon-Nikodym derivative/density w.r.t.\ $P^Y$ of the following measure on $\Ycal$: 
     \[ E_\pm(B) := \E\lB X_\pm \cdot \I_B(Y) \rB = \int x \cdot \I_B(y) \, dP^{(X_\pm,Y)}(x,y).  \]
     One can easily see that $E_\pm \ll P^Y$ and that the densities exist by the Radon-Nikodym theorem.
% \item This defines a measurable map:
 %    \[   \Ycal \to \bar\R,\quad y \mapsto \E[X|Y=y].\]
     %This means that $y \mapsto \E[X|Y=y]$ is uniquely defined via:
     %\[ \E\lB X \cdot \I_B(Y)\rB = \int_B \E[X|Y=y]\,P^Y(dy),   \]
    %for all measurable $B \ins \Ycal$.
 \item The \emph{conditional expectation}  of $X$ given $Y$ is then the measurable map defined via:
     \[\E[X| Y] :\, \Wcal \to \bar\R,\quad w \mapsto \E[X|Y](w):=\E[X|Y=Y(w)]=\E[X|Y=y]|_{y=Y(w)},\]
     i.e.\ the composition of $Y$ with the measurable map $y \mapsto\E[X|Y=y]$.
     %\[ \E[X|\ul Y] = \E[X|Y] \circ Y.  \]
\end{enumerate}    
\end{Def}

\begin{Rem}
    The construction from above also works with a measure $\mu$ such that $\mu^Y$ is $\sigma$-finite (instead of $P$) since we only need to guarantee the existence of the Radon-Nikodym derivative.
\end{Rem}

\begin{Not}
 Let $\Wcal$, $\Zcal$, $\Ycal$ be measurable spaces and $Z:\, \Wcal \to \Zcal$ and $Y:\,\Wcal \to \Ycal$ be measurable maps.
 We write:
 \[ Z \precsim Y  \]%\qquad \text{ or } \qquad Z \in \sigma(Y) \]% \qquad \text{ or } \qquad \sigma(Z) \ins \sigma(Y) \]
 if there exists a measurable function $F: \Ycal \to \Zcal$ % 
 %$F:Y(\Wcal) \to \Zcal$ 
 such that $Z= F \circ Y $; in other words if $Z$ is a deterministic (measurable) function of $Y$, i.e.: $Z=F(Y)$.\\
 If $\mu$ is a measure on $\Wcal$ we also write:
 \[ Z \precsim_\mu Y\]
 if there exists a measurable map $F$ such that $Z=F(Y)$ $\mu$-almost-surely.
 %if there exists a measurable $\mu$-null set $N \ins \Wcal$ and a measurable map $F:Y(N^\cmpl) \to \Zcal$ such that for all $w \in N^\cmpl$ we have $Z(w)=F(Y(w))$.

\end{Not}

\begin{Thm}
 Let $(\Wcal,P)$ be a probability space
 and $X,T:\, \Wcal \to \R$, $Y:\,\Wcal \to \Ycal$, $Z:\, \Wcal \to \Zcal$ be random variables with $\E[|X|]< \infty$ (or as long as we do not run into the ``$\infty - \infty=?$'' problem).
    Then we have the following properties:
    \begin{enumerate}
        \item $\E[X|Y]$ is the unique real valued random variable $Z$ (up to $P$-null set)  such that: 
            \begin{enumerate}
                \item $Z \precsim_P Y$ and:
                \item 
            for all measurable $B \ins \Ycal$:            
            \[ \E\lB Z \cdot \I_B(Y) \rB = \E[ X \cdot \I_B(Y)].  \]
            \end{enumerate}
%        \item  For all  measurable  $g:\,\Ycal \to \R$ with $\E[|g(Y) \cdot X|]< \infty$ we have:
%            \[ \E\lB g(Y) \cdot X|Y \rB = g(Y) \cdot \E[X|Y] \quad P\text{-a.s.}  \]
        \item  For all real valued random variables $Z \precsim_P Y$  with $\E[|Z \cdot X|]< \infty$ we have:
            \[ \E\lB Z \cdot X|Y \rB = Z \cdot \E[X|Y] \quad P\text{-a.s.}  \]
        \item Linearity: For all $a,b \in \R$ we have:
            \[ \E[ a \cdot X + b\cdot T|Y] = a \cdot \E[X|Y] + b\cdot\E[T|Y] \quad P\text{-a.s.}  \]
        \item Constants: $\E[1|Y] = 1$ $P$-a.s.
        \item Constant maps: If $Y$ is a constant map then: $\E[X|Y] = \E[X]$ $P$-a.s.
        \item Independence (see \ref{ci-random-variables}): If $X \Indep Y$ then: $\E[X|Y] = \E[X]$ $P$-a.s.    
        \item Deterministic dependence: If $X \precsim_P Y$ then: $\E[X|Y] = X$ $P$-a.s.   
        \item Monotonicity: If $ X \ge T$ $P$-a.s. then we have:
            \[ \E[X|Y] \ge  \E[T|Y] \quad P\text{-a.s.}    \]
        \item Jensen inequality: Let $\varphi: \R \to \R$ be convex then we have:
            \[ \varphi\lp \E[X|Y]\rp \le \E[ \varphi(X)|Y] \quad P\text{-a.s.}   \]
        \item Triangle inequality: $|\E[X|Y]| \le \E[\,|X|\,|\, Y]$ $P$-a.s.
%        \item Tower rule: For all measurable $g:\, \Ycal \to \Zcal$ we have:
%            \[ \E\lB\E[X|Y]|g(Y)\rB = \E\lB\E[X|g(Y)]|Y\rB = \E[X|g(Y)]\quad P\text{-a.s.}    \]
        \item Tower rule: If $Y \precsim Z$ then:
            \[ \E\lB\E[X|Y]|Z\rB = \E\lB\E[X|Z]|Y\rB = \E[X|Y]\quad P\text{-a.s.}    \]
        \item Tower rule, special case:
            \[ \E\lB\E[X|Y]|Y,Z\rB = \E\lB\E[X|Y,Z]|Y\rB = \E[X|Y]\quad P\text{-a.s.}    \]
        \item Monotone convergence, dominated convergence, etc. (see literature).
    \end{enumerate}
\end{Thm}


\subsection{The Lebesgue Measure}

%Remark: continous mappings are measurable wrt Borel sigma algebra
% Existence of Lebegue measure and change of density formula

\begin{comment}
\begin{Thm}[Lebesgue measure]
    \label{thm-lebesgue-measure}
    Consider $\R^D$ with its Borel $\sigma$-algebra $\Bcal_{\R^D}$.
    Then there exists a unique measure $\lambda^D$ on $(\R^D,\Bcal_{\R^D})$ such that:
    \[ \lambda^D\lp [a_1,b_1] \times \cdots \times [a_D, b_D]\rp = (b_1-a_1) \cdots (b_D-a_D)   \]
    for all $a_d, b_d \in \R$, $a_d \le b_d$, $d=1,\dots,D$. 
    It will be called the $D$-dimensional \emph{Lebesgue measure}.
    If the dimension is clear from the context we might just write $\lambda$ for $\lambda^D$.
\end{Thm}
\end{comment}

\begin{Def}[The Lebesgue (outer) measure]
    \label{thm-lebesgue-measure}
    The \emph{Lebesgue (outer) measure} $\lambda^D$ on $\R^D$ is given for subsets $A \ins \R^D$ via:
    \[ \lambda^D(A) := \inf\lC \sum_{n \in \N} \mathrm{vol}^D\lp [a^{(n)},b^{(n)}] \rp  \st  
    A \ins \bigcup_{n \in \N} [a^{(n)},b^{(n)}]  \rC, \]
    where the infimum is running over sequences of $D$-dimensional cubes:
     \[ [a^{(n)},b^{(n)}] = [a_1^{(n)},b_1^{(n)}] \times \cdots \times [a_D^{(n)}, b_D^{(n)}],\]
  with $a^{(n)}=(a_1^{(n)},\dots,a_D^{(n)})$, $b^{(n)}=(b_1^{(n)},\dots,b_D^{(n)}) \in \R^D$, $a_d^{(n)} \le b_d^{(n)}$ for $d=1,\dots,D$, $n \in \N$, that jointly cover $A$, where the $D$-dimensional volume is given by:
  \[ \mathrm{vol}^D\lp [a^{(n)},b^{(n)}] \rp := (b_1^{(n)}-a_1^{(n)}) \cdots (b_D^{(n)}-a_D^{(n)}),\qquad \mathrm{vol}^D\lp \emptyset \rp :=0. \]
\end{Def}

\begin{Thm}[The Lebesgue measure]
 The \emph{Lebesgue measure} $\lambda^D$, when restricted to the Borel-$\sigma$-algebra of $\R^D$,
  is the unique measure on $\R^D$ that satisfies:
    \[ \lambda^D\lp [a,b] \rp = \mathrm{vol}^D\lp [a,b] \rp,     \]
    for all $D$-dimensional cubes $[a,b]$.
    If the dimension is clear from the context we might just write $\lambda$ for $\lambda^D$.
\end{Thm}

\begin{Thm} Let $\lambda$ be the Lebesgue measure on the interval $[a,b]$.  Let $f:\, [a,b] \to \R$ be a Riemann integrable function (e.g.\ a continuous function) then $f$ is also $\lambda$-integrable and we have:
    \[ \int_a^b f(x)\,dx = \int_{[a,b]}  f(x)\, \lambda(dx). \]
    %Roughly, this means that we can use all the rules we know about integration to compute Lebesgue integrals.
\end{Thm}

\begin{Thm}[Fundamental theorem of calculus]
   % \begin{enumerate}
    %    \item 
    Let $f: \R \to \R$ be a measurable function such that $\int_{[a,b]}  |f| \,d\lambda < \infty$ for $a,b \in \R$. % all compact subsets $K \ins \R$.
            For fixed $c \in \R$ define $F:\, [c,\infty) \to \R$ via:
            \[ F(x) := \int_{(c,x]} f\,d\lambda.  \]
            Then $F$ is differentiable in  $\lambda$-almost-all $x \in \R$ and for those points we have:
            \[ F'(x) = f(x).\]
    %    \item
    %\end{enumerate}
\end{Thm}

\subsection{Transformation Rules}


\begin{Thm}[General integral transformation]
    Let $(\Wcal,\mu)$ be a measure space and $X :\,\Wcal \to \Xcal$ and $F:\, \Xcal \to \R$ be measurable.
    Then we have:
    \[ \int F(X)\,d\mu = \int F\, d(X_*\mu),  \]
    if either side is well-defined.
    Written in longer form this is:
    \[ \int F(X(w))\,\mu(dw) = \int F(x)\, (X_*\mu)(dx).  \]
\end{Thm}



\begin{Thm}[Push-forward of densities]
    Let $(\Wcal,\mu)$ be a  measure space and $\nu$ another measure on $\Wcal$. Let $\varphi:\,\Wcal \to \Ycal$ be a measurable mapping such that 
    $\varphi_*\mu$ is $\sigma$-finite.
    If $\nu$ has a density w.r.t.\ $\mu$ then the push-forward measure $\varphi_*\nu$ has a density w.r.t.\ $\varphi_*\mu$ given as follows:
    %For this, note that $\varphi_*\nu$ can be written in terms of the density $\frac{d\nu}{d\mu}$ as follows:
    %\[ (\varphi_*\nu) = \int \I_{\varphi^{-1}(B)}(w)\cdot \frac{d\nu}{d\mu}(w)\,d\mu(w). \]
    \[ \frac{d(\varphi_*\nu)}{d(\varphi_*\mu)}(y) = \E_\mu\lB \frac{d\nu}{d\mu} \bigg| \varphi=y\rB  = "\int \frac{d\nu}{d\mu}(w) \, \mu(dw|\varphi=y)", \]
    for $\varphi_*\mu$-almost-all $y \in \Ycal$,  where the conditional integral $\E_\mu$  is constructed 
    the same way as the conditional expectation but using the $\sigma$-finite measure $\varphi_*\mu$.\\
    %the conditional measure might not always exists, but the conditional integral as whole does!!!!
    If, furthermore, $\varphi$ is a measurable isomorphism then we get: 
    \[ \frac{d(\varphi_*\nu)}{d(\varphi_*\mu)}(y) = \frac{d\nu}{d\mu}(\varphi^{-1}(y)) \]
     for $\varphi_*\mu$-almost-all $y \in \Ycal$.
\end{Thm}

\begin{Thm}[Transformation formula for the Lebesgues measure]
    \label{transformation-formula}
    Let $\varphi:\, \R^D \to \R^D$ be a continously differentiable bijection of $\R^D$ (or of open/closed subsets therein) with Jacobian $\varphi'(x)$ at point $x$. 
    Let $\lambda$ be the Lebesgue measure on $\R^D$.
    Then $\varphi_* \lambda$ is absolute continuous w.r.t.\ $\lambda$ with density given by:
    \[ \frac{d (\varphi_*\lambda)}{d\lambda}(y) = | \det \varphi' (\varphi^{-1}(y))|^{-1}  \]
    for all $y \in \R^D$ (or in that open/closed subset, and $=0$ outside).
\end{Thm}


\begin{Cor}[Transformation of (probability) densities w.r.t.\ the Lebesgue measure]\label{lebesgue-density-transformation}
    Let the setting be like in \ref{transformation-formula}.
    Let $\nu$ be a (probability) measure on $\R^D$ with (probability) density $p$ w.r.t.\ $\lambda$.
    Then $\varphi_*\nu$ also has a (probability) density w.r.t.\ $\lambda$, which is then given by:
    \[ \frac{d(\varphi_*\nu)}{d\lambda}(y) = \frac{d(\varphi_*\nu)}{d(\varphi_*\lambda)}(y) \cdot \frac{d(\varphi_*\lambda)}{d\lambda}(y)  =  p(\varphi^{-1}(y)) \cdot |\det \varphi'(\varphi^{-1}(y)|^{-1}.   \]
\end{Cor}

\begin{Thm}[A bit more general, \cite{Fremlin} Cor.\ 263F, 262F(b)]
    Let $\Xcal \ins \R^D$ be a measurable set 
    and $\varphi:\, \Xcal \to \R^D$ an injective Lipschitz function.
Let $\Xcal' \ins \Xcal$ be the set of points $x$ at which $\varphi$ has a derivative $\varphi'(x)$ relative to $\Xcal$\footnote{We say that $\varphi$ is \emph{differentiable relative to $\Xcal$} at $x \in \Xcal$  if there exists  $\varphi'(x) \in \R^{D \times D}$ such that for every $\epsilon >0$ there exists a $\delta >0$ such that for all
    $y \in \Xcal$ with $\|y-x\| < \delta$  we have that: 
$\|\varphi(y)-\varphi(x) - \varphi'(x) \cdot (y-x) \| \le \epsilon \cdot \|y-x\| $. Note that in this definition such a derivative $\varphi'(x)$ does not need to be unique.}.
Then we have:
    \begin{enumerate}
        \item $\Xcal \sm \Xcal'$ is a $\lambda$-null set.
        \item $|\det \varphi'|:\, \Xcal' \to [0,\infty)$ is measurable.
        \item $\varphi(\Xcal) \ins \R^D$ is a measurable set.
        \item $\lambda(\varphi(\Xcal)) = \int_\Xcal |\det \varphi'(x)| \,d\lambda(x)$.
        \item For every real-valued function $g$ defined on a subset $\Ycal \ins \varphi(\Xcal)$
            we have:
            \[ \int_{\varphi(\Xcal)} g(y)\,d\lambda(y) = \int_\Xcal g(\varphi(x)) \cdot |\det \varphi'(x)| \, d\lambda(x),  \]
            if either integral is defined in $[-\infty,\infty]$ and provided we interpret
            $g(\varphi(x)) \cdot |\det \varphi'(x)| := 0$ if $\varphi(x) \notin \Ycal$ and $|\det \varphi'(x)| =0$. 
    \end{enumerate}
\end{Thm}


\begin{Rem}[Transformation rule for discrete measures]
    Let $\Xcal$ be a measurable space and $\mu$ be a discrete (probability) measure on $\Xcal$ supported on the countable discrete subset $\Wcal \ins \Xcal$ with mass function given by:
    \[ m(x) = \frac{d\mu}{d\#_\Wcal}(x),\]
    where $\#_\Wcal$ is the counting measure w.r.t.\ $\Wcal$ given by: $\#_\Wcal(A) := \#(\Wcal \cap A)$.
    Let $\varphi:\,\Xcal \to \Ycal$ be a measurable map. Then $\varphi_*\mu$ is a discrete measure supported on $\varphi(\Wcal)$ 
    with mass function/density:
    \[ \frac{d\varphi_*\mu}{d\#_{\varphi(\Wcal)}}(y) = \sum_{ w \in \varphi^{-1}(y) \cap \Wcal } m(w).
    \]
\end{Rem}

\begin{Eg}[Linear transformation of Gaussian distributions]
\end{Eg}

\begin{Eg}[Density of Chi-square distributions]
\end{Eg}





\begin{comment}
\subsection{Exercises}


\begin{Exc}[$\sigma$-algebra]
    \label{exercise-sigma-algebra}
    \begin{enumerate}
        \item
    Let $\Wcal$ be a set and $\Bcal$ a $\sigma$-algebra on $\Wcal$. Show the following:
    \begin{enumerate}
        \item If $A_1,A_2 \in \Bcal$ then also: $A_1 \cup A_2 \in \Bcal$.
        \item If $A_n \in \Bcal$ for $n \in \N$ then also: $\bigcap_{n \in \N} A_n \in \Bcal$.
        \item If $A_1,A_2 \in \Bcal$ then also: $A_1 \sm A_2 \in \Bcal$.
    \end{enumerate}
\item Let $\Wcal$ be a set. Show that $2^\Wcal$ is a $\sigma$-algebra on $\Wcal$.
\item Let $\Wcal$ be a set and $(\Bcal_i)_{i \in I}$ be an arbitrary family of $\sigma$-algebras of $\Wcal$. Show that the intersection:
    $  \bigcap_{i \in I} \Bcal_i $
    is again a $\sigma$-algebra on $\Wcal$.
    \end{enumerate}
\end{Exc}


\begin{Exc}[Borel $\sigma$-algebra]
    \label{exercise-Borel-sigma-algebra}
    \begin{enumerate}
        \item 
    Show that the Borel $\sigma$-algebra $\Bcal_\R$ on $\R$ can also be generated by the following sets of subsets of $\R$:
    \begin{enumerate}
        \item $\lC (-\infty,b]\;|\; b \in \Q   \rC$,
       % \item  $\lC (a,b]\;|\; a,b \in \R, a\le b \rC$,
        \item  $\lC (a,b)\;|\; a,b \in \R, a\le b \rC$,
        \item  $\lC [a,b]\;|\; a,b \in \R, a\le b \rC$,
        \item etc.
    \end{enumerate}
\item Show that the Borel $\sigma$-algebra $\Bcal_{\R^D}$ is also the product $\sigma$-algebra 
    of the 1-dimensional Borel $\sigma$-algebras $\Bcal_\R$, i.e.\ show:
    \[ \Bcal_{\R^D} = \Bcal_\R \otimes \cdots \otimes \Bcal_\R.\]
%\[ \Bcal_{\R^D} = \sigma\lp\lC (-\infty,b_1] \times \cdots \times (-\infty,b_D]\;|\; b_1,\dots,b_D \in \R    \rC\rp.  \]
    \end{enumerate}
\end{Exc}


\begin{Exc}[$\sigma$-algebras induced by mappings]
    \begin{enumerate}
        \item Show that the composition $g \circ f$ of two measurable mappings $f: \Wcal \to \Zcal$ and $g: \Zcal \to \Ycal$ is again measurable w.r.t.\ the corresponding $\sigma$-algebras.
        \item Let $f: \Wcal \to \Zcal$ be any mapping and $\Bcal_\Zcal$ be a $\sigma$-algebra on $\Zcal$. Then show that:
            \[ f^*\Bcal_\Zcal := \lC f^{-1}(C)\,|\, C \in \Bcal_\Zcal  \rC\]
            is a $\sigma$-algebra on $\Wcal$. Furthermore, show that $f$ is $f^*\Bcal_\Zcal$-$\Bcal_\Zcal$-measurable.
            Show that $f^*\Bcal_\Zcal$ is the smallest $\sigma$-algebra $\Bcal_\Wcal$ on $\Wcal$ that makes $f$ $\Bcal_\Wcal$-$\Bcal_\Zcal$-measurable.
        \item Let $f: \Wcal \to \Zcal$ be any mapping and $\Bcal_\Wcal$ be a $\sigma$-algebra on $\Wcal$. Then show that:
            \[ f_*\Bcal_\Wcal := \{ C \ins \Zcal\,|\, f^{-1}(C) \in \Bcal_\Wcal  \}\]
            is a $\sigma$-algebra on $\Zcal$. Furthermore, show that $f$ is $\Bcal_\Wcal$-$f_*\Bcal_\Wcal$-measurable.
            Shwo that $f_*\Bcal_\Wcal$ is the biggest $\sigma$-algebra $\Bcal_\Zcal$ on $\Zcal$ that makes $f$ $\Bcal_\Wcal$-$\Bcal_\Zcal$-measurable.
         \item Let $f: \Wcal \to \Zcal$ be any mapping and $\Bcal_\Wcal$ and $\Bcal_\Zcal$ be $\sigma$-algebras on $\Wcal$ and $\Zcal$, resp..
             Show that $f$ is measurable iff $f^*\Bcal_\Zcal \ins \Bcal_\Wcal$ iff $\Bcal_\Zcal \ins f_*\Bcal_\Wcal$.
    \end{enumerate}
\end{Exc}

\begin{Exc}
    Show that the subspace $\sigma$-algebra $\Bcal_{|\Zcal}$ really is a $\sigma$-algebra on $\Zcal$.
\end{Exc}

\begin{Exc}
    Show that the push-forward measure really is a measure.
\end{Exc}


\begin{Exc}
    Give a more explicit description of the product $\sigma$-algebra in terms of product sets $A_1 \times \cdots \times A_N$.
\end{Exc}
\end{comment}


\subsection{Measure Extension Theorems}

\begin{Thm}[Measure extension theorem, see \cite{Kle20} Thm.\ 1.53, Thm.\ 1.36]
    \label{thm:ext-thm-2}
    Let $\Acal$ be a ring (e.g.\ an algebra) of subsets of a set $\Xcal$.
    Let $\mu:\, \Acal \to [0,\infty)$ be a (finitely) additive set function 
    with $\mu(\emptyset) =0$ that is also
    \emph{$\emptyset$-continuous}:
    \begin{align}
        \inf_{n \in \N} \mu(A_n) = 0, 
    \end{align}
    for all non-increasing sequences $(A_n)_{n \in \N}$ with $A_n \in \Acal$ and $A_{n+1} \ins A_n$ for all $n \in \N$ and $\bigcap_{n \in \N} A_n = \emptyset$.

    Then there exists a unique $\sigma$-finite measure $\nu:\, \sigma(\Acal) \to [0,\infty]$ such that $\nu(A) = \mu(A)$ for all $A \in \Acal$.
\end{Thm}



\begin{Thm}[Ionescu-Tulcea extension theorem, see \cite{Tul49,Lam87}]
    \label{thm:ionescu-tulcea}
    Let $I$ be an arbitrary set (not necessarily countable) and $(\Xcal_i,\Bcal_i)$, $i \in I$, measurable spaces.
    For subsets $J \ins I$ we put:
    \begin{align}
        \Xcal_J &:= \prod_{j \in J} \Xcal_j, & \Bcal_J&:= \bigotimes_{j \in J} \Bcal_j,
    \end{align}
    the product space endowed with its product $\sigma$-algebra.
    Now assume that we have a probability measure $\mu_J$ on $(\Xcal_J,\Bcal_J)$ 
    for every finite subset $J \ins I$ such that:
    \begin{enumerate}
        \item for every finite subsets $L \ins J \ins I$ we have:
            $\pr_{L,*}\mu_J = \mu_L$,
        \item for every finite subset $J \ins I$ and $i \in I \sm J$ there exists a Markov kernel: $\mu_{i|J}:\, \Xcal_J \dshto \Xcal_i$ such that:
            \begin{align}
                \mu_{\{ i \} \dcup J} &= \mu_{i|J} \otimes \mu_J.
            \end{align}
    \end{enumerate}
    Then there exists a probability measure $\mu_I$ on $(\Xcal_I,\Bcal_I)$ such that
    for every finite subset $J \ins I$ we have:
            \begin{align}
                \pr_{J,*}\mu_I = \mu_J.
            \end{align}
\begin{proof}
    We first put, with $\pr_J:\, \Xcal_I \to \Xcal_J$ the canonical projections:
    \begin{align}
        \Acal &:= \bigcup_{\substack{J \ins I\\ \#J < \infty}} \pr_J^*\Bcal_J
        = \lC \pr_J^{-1}(B) \ins \Xcal_I \st J \ins I, \#J < \infty, B \in \Bcal_J \rC.
    \end{align}
    Then, per definition, $\Bcal_I = \sigma(\Acal)$.
    Furthermore, $\Acal$ is an algebra of subsets of $\Xcal_I$.
    Indeed, let $A_1,A_2 \in \Acal$ then $A_l \in \pr_{J_l}^{*}\Bcal_{J_l}$ for some finite subsets $J_l \ins I$, $l=1,2$. Then $J:=J_1 \cup J_2$ is also a finite subset of $I$ and we have: $A_1, A_2 \in \pr_J^{*}\Bcal_J$. So, 
    $A_l = \pr_J^{-1}(B_l)$ for some $B_l \in \Bcal_J$, $l=1,2$.
    This then shows that:
    \begin{align}
        A_l^\cmpl &= \pr_J^{-1}(B_l^\cmpl) &\in \pr_J^{*}\Bcal_J \ins \Acal, \\
        A_1 \cup A_2 &= \pr_J^{-1}(B_1 \cup B_2) &\in \pr_J^{*}\Bcal_J \ins \Acal, \\
        A_1 \cap A_2 &= \pr_J^{-1}(B_1 \cap B_2) &\in \pr_J^{*}\Bcal_J \ins \Acal.
    \end{align}
    It is also clear that: $\Xcal_I,\emptyset \in \Acal$. 
    So, $\Acal$ is an algebra of subsets of $\Xcal_I$.

    We can now define the set function $\mu:\, \Acal \to [0,1]$ via:
    \begin{align}
        \mu(A) &:= \mu_J(B),  \quad\text{ for }\quad A =  \pr_J^{-1}(B), \quad B \in \Bcal_J.
    \end{align}
    This is well-defined because of the condition: $\pr_{L,*}\mu_J = \mu_L$ for finite subsets $L \ins J \ins I$.
    
    It is also clear that $\mu$ is additive. Indeed, if $A_1,A_2 \in \Acal$ are disjoint then $A_l = \pr_J^{-1}(B_l)$ for some finite subset $J \ins I$ and some disjoint $B_l \in \Bcal_J$, $l=1,2$. The additivity of $\mu_J$ then shows the additivity of $\mu$:
    \begin{align}
        \mu(A_1 \dcup A_2) = \mu_J(B_1 \dcup B_2) = \mu_J(B_1) + \mu_J(B_2) 
        = \mu(A_1) + \mu(A_2).
    \end{align}
    To apply the extension theorem \ref{thm:ext-thm-2} it is left to check that $\mu$ 
    is $\emptyset$-continuous on $\Acal$.
    For this, and, by way of contradiction, consider a non-increasing sequence $A_n \in \Acal$, $n \in \N$, with $\bigcap_{ n \in \N} A_n = \emptyset$ and $\inf_{n \in \N} \mu(A_n) > \epsilon > 0$. We can assume that $A_n = \pr_{J_n}^{-1}(B_n)$ with $B_n \in \Bcal_{J_n}$ with the inclusion of finite subsets:
    $J_n \ins J_{n+1} \ins I$ for all $n \in \N$. We totally order the countable set $\bigcup_{n \in \N} J_n$ such that $k < l $ if $k \in J_n$ and $l \in J_{n+1} \sm J_n$.

We introduce the following abbreviations for $n \in \N$:
\begin{align}
    \Ycal_n &:= \Xcal_{J_n \sm J_{n-1}}, & \Ycal_{\le n} &:= \prod_{l=1}^n Y_n, &
    \Ycal_\cmpl &:= \Xcal_{I \sm \bigcup_{n \in \N} J_n},\\
    \mu_{n|<n} &:= \bigotimes_{k \in J_n \sm J_{n-1}} \mu_{k|\lC l \in J_n \st l <k\rC \cup J_{n-1}}, &     \mu_{\le n} &:= \mu_{J_n}. 
\end{align}
    Then $\Ycal := \Ycal_\cmpl \times \prod_{n \in \N} \Ycal_n = \Xcal_I$.
   We also put:
\begin{align}
    h_n(y) &:=  g_n(y_{\le n}) := \I_{B_n}(y_{\le n}) = \I_{A_n}(y), 
    & h(y) &:= \inf_{n \in \N} h_n(y).
\end{align}
    By assumption, $\bigcap_{n \in \N} A_n = \emptyset$, we have that $h(y) =0 $ for all $y \in \Ycal$.
    Since $A_n \ins A_{n-1}$ we have for all $n \in \N$:
    \begin{align}
        0 = h(y) \le   h_n(y) \le h_{n-1}(y) \le 1.
    \end{align}
We define for $k, n \in \N$:
\begin{align}
    f_n^{(k)}(y_{\le k}) &:= \int g_n(y_{k+1:n},y_{\le k})\, \mu_{k+1:n|\le k}(dy_{k+1:n}|y_{\le k}),\\
    f^{(k)}(y_{\le k}) &:= \inf_{n \in \N} f_n^{(k)}(y_{\le k}).
\end{align}
Note that we then also have for all $k, n \in \N$ and $y \in \N$:
    \begin{align}
        0 \le f^{(k)}(y_{\le k}) \le  f_n^{(k)}(y_{\le k})  \le f_{n-1}^{(k)}(y_{\le k})  \le 1.
    \end{align}
We also put for $k \in \N$:
    \begin{align}
        C_{\le k} &:= \lC y_{\le k} \in \Ycal_{\le k} \st f^{(k)}(y_{\le k}) > \epsilon \rC.
    \end{align}

    By the above assumption we have for every $n \in \N$:
    \begin{align}
        0 < \epsilon &< \inf_{n \in \N} \mu(A_n) \\
        &= \inf_{n \in \N}  \mu_{\le n}(B_n)\\
        &=\inf_{n \in \N} \E_{\mu_{\le n}}[g_n] \\
        &=\inf_{n \in \N} \int \int g_n(y_{2:n},y_1)  \, \mu_{2:n|1}(dy_{2:n}|y_1) \, \mu_1(dy_1)\\
        &=\inf_{n \in \N} \int f_n^{(1)}(y_1) \, \mu_1(dy_1)\\
        &= \int \inf_{n \in \N} f_n^{(1)}(y_1) \, \mu_1(dy_1)\\
        &= \int f^{(1)}(y_1) \, \mu_1(dy_1).
    \end{align}
Where the integral and infimum can be interchanged because $f^{(1)}$ is $\Bcal_1$-measurable and the monotone convergence theorem, applied to $1-f_n^{(1)}$.
    We see that: $\mu_1(C_1) > 0$. Otherwise, $\int f^{(1)}(y_1) \, \mu_1(dy_1) \le \epsilon$, which would contradict the above sequence of inequalities.

 Now inductively for $k \in \N$ and $y_{\le k} \in C_{\le k}$ we have:
   \begin{align}
   0 < \epsilon &< f^{(k)}(y_{\le k}) \\
                &=\inf_{n \in \N}  f_n^{(k)}(y_{\le k}) \\
                &=\inf_{n \in \N} \int f_n^{(k+1)}(y_{k+1},y_{\le k}) \, \mu_{k+1|\le k}(dy_{k+1}|y_{\le k})\\
                &= \int \inf_{n \in \N} f_n^{(k+1)}(y_{k+1},y_{\le k}) \, \mu_{k+1|\le k}(dy_{k+1}|y_{\le k})\\
                &= \int f^{(k+1)}(y_{k+1},y_{\le k}) \, \mu_{k+1|\le k}(dy_{k+1}|y_{\le k}).
    \end{align}
    This shows that: $\mu_{k+1|\le k}(C_{\le k+1}^{y_{\le k}}|y_{\le k}) > 0$ 
    for $y_{\le k} \in C_{\le k}$.
    This means that we can inductively construct a $y \in \Ycal$ with components: 
    $y_1 \in C_1$ and $y_{k+1} \in C_{\le k+1}^{y_{\le k}}$ for $k \in \N$, 
    and an arbitrary $y_\cmpl \in \Ycal_\cmpl$. 
    This $y$ then satisfies $h_n(y) > \epsilon > 0$ for all $n \in \N$ and thus 
    $h(y) = \inf_{n \in \N} h_n(y) \ge \epsilon > 0$, 
    which lies in contradiction to $h(y) =0$ for all $y \in \Ycal$.

    This shows that $\mu$ is $\emptyset$-continuous.
    It follows by the extension theorem \ref{thm:ext-thm-2} that $\mu$ has a unique extension to a probability measure to $\sigma(\Acal)=\Bcal_I$.
    This shows the claim.
\end{proof}
\end{Thm}

\begin{Cor}[Ionescu-Tulcea extension theorem for Markov kernels]
    \label{cor:ionescu-tulcea-markov-kernels}
    Let $I$ be an arbitrary set (not necessarily countable) and $(\Zcal,\Bcal_\Zcal)$ and $(\Xcal_i,\Bcal_i)$, $i \in I$, measurable spaces.
    For subsets $J \ins I$ we put:
    \begin{align}
        \Xcal_J &:= \prod_{j \in J} \Xcal_j, & \Bcal_J&:= \bigotimes_{j \in J} \Bcal_j,
    \end{align}
    the product space endowed with its product $\sigma$-algebra.
    Now assume that for every finite subset $J \ins I$  we are given a Markov kernel:
    \begin{align*}
        K_J(X_J|Z):\, \Zcal \dshto \Xcal_J,
    \end{align*}
     such that:
    \begin{enumerate}
        \item for every finite subsets $L \ins J \ins I$ we have:
          \begin{align*}
              K_J(X_L|Z) = K_L(X_L|Z):\, \Zcal \dshto \Xcal_L,
          \end{align*}
        \item for every finite subset $J \ins I$ and $i \in I \sm J$ there exists a Markov kernel:
          \begin{align*}
              K_{i|J}(X_i|X_J,Z):\, \Xcal_J\times\Zcal \dshto \Xcal_{\{i\} \dcup J},
          \end{align*}
           such that:
          \begin{align*}
              K_{i|J}(X_i|X_J,Z) \otimes K_J(X_J|Z) & = K_{\{i\} \dcup J}(X_i,X_J|Z)  :\, \Zcal \dshto \Xcal_{\{i\} \dcup J}.
          \end{align*}
    \end{enumerate}
    Then there exists a Markov kernel:
    \begin{align*}
        K(X_I|Z):\, \Zcal \dshto \Xcal_I,
    \end{align*}
    such that for every finite subset $J \ins I$ we have:
    \begin{align*}
        K(X_J|Z) &= K_J(X_J|Z):\, \Zcal \dshto \Xcal_J.
    \end{align*}
\begin{proof}
    For every $z \in \Zcal$ we can apply the Ionescu-Tulcea extension theorem \ref{thm:ionescu-tulcea} separately and get a probability measure 
    $K(X_I|Z=z)$ on $\Bcal_I$ such that for every finite subset $J \ins I$ we have:
    \begin{align*}
        K(X_J|Z=z) &= K_J(X_J|Z=z).
    \end{align*}
    We are left to check that the map:
    \begin{align*}
        K(X_I|Z):\, \Zcal \to \Prob(\Xcal_I), \qquad z \mapsto K(X_I|Z=z),
    \end{align*}
    is measurable. By Dynkin's lemma and the definition of the product $\sigma$-algebra $\Bcal_I$ this only needs to be checked on sets 
    $A_J \in \Bcal_J$ for finite subsets $J \ins I$. Since we have:
    \begin{align*}
        K(X_I \in \pr_{J}^{-1}(A)|Z=z) &= K(X_J \in A|Z=z) = K_J(X_J \in A|Z=z),
    \end{align*}
    and $z \mapsto K_J(X_J \in A|Z=z)$ is measurable for finite $J \ins I$, the claim follows.
\end{proof}
\end{Cor}


\begin{Cor}[Kolmogorov extension theorem for Markov kernels]
    \label{cor:kolmogorov}
    Let $I$ be an arbitrary set (not necessarily countable) $(\Zcal,\Bcal_\Zcal)$ a measurable space and $(\Xcal_i,\Bcal_i)$, $i \in I$, 
    \emph{standard} measurable spaces.
    For subsets $J \ins I$ we put:
    \begin{align}
        \Xcal_J &:= \prod_{j \in J} \Xcal_j, & \Bcal_J&:= \bigotimes_{j \in J} \Bcal_j,
    \end{align}
    the product space endowed with its product $\sigma$-algebra.
    Now assume that for every finite subset $J \ins I$  we are given a Markov kernel:
    \begin{align*}
        K_J(X_J|Z):\, \Zcal \dshto \Xcal_J,
    \end{align*}
     such that for every finite subsets $L \ins J \ins I$ we have:
          \begin{align*}
              K_J(X_L|Z) = K_L(X_L|Z):\, \Zcal \dshto \Xcal_L.
          \end{align*}
    Then there exists a Markov kernel:
    \begin{align*}
        K(X_I|Z):\, \Zcal \dshto \Xcal_I,
    \end{align*}
    such that for every finite subset $J \ins I$ we have:
    \begin{align*}
        K(X_J|Z) &= K_J(X_J|Z):\, \Zcal \dshto \Xcal_J.
    \end{align*}
\begin{proof}
    This directly follows from Ionescu-Tulcea extension theorem for Markov kernels \ref{cor:ionescu-tulcea-markov-kernels} and the fact that on standard measurable spaces we always have conditional Markov kernels by the disintegration theorem \ref{thm-regular-conditional-Markov-kernel}.
\end{proof}
\end{Cor}
